\section{The content projector returns the full shell}
\label{sec:content-projector}

The small-content indicator in \eqref{eq:literal-AC} admits an exact
common-divisor M\"obius expansion.  At nonzero determinant frequency,
this expansion can expose useful separated phases.  At \(r=0\), its
principal term is the original unrestricted amplitude.

\subsection{Cumulative common-divisor inversion}

For \(C\ge1\), define
\begin{equation}
 \eta_{\le C}(t)
 =\sum_{\substack{c\mid t\\c\le C}}\mu(t/c).
 \label{eq:content-eta}
\end{equation}

\begin{lemma}[Exact cumulative projector]
\label{lem:content-projector}
For every positive integer \(G\),
\begin{equation}
 \boxed{
 \one_{G\le C}
 =\sum_{t\mid G}\eta_{\le C}(t).}
 \label{eq:content-projector}
\end{equation}
Moreover,
\begin{equation}
 \eta_{\le C}(1)=1,
 \qquad
 \boxed{
 \sum_{\substack{t\mid G\\t>1}}\eta_{\le C}(t)
 =-\one_{G>C}.}
 \label{eq:content-nontrivial-reassembly}
\end{equation}
\end{lemma}

\begin{proof}
Reversing the divisor sums gives
\begin{align*}
 \sum_{t\mid G}\eta_{\le C}(t)
 &=
 \sum_{\substack{c\mid G\\c\le C}}
 \sum_{k\mid G/c}\mu(k).
\end{align*}
The inner sum equals one precisely when \(c=G\) and is zero
otherwise.  This proves \eqref{eq:content-projector}.  The value at
\(t=1\) is immediate from \eqref{eq:content-eta}.  Subtracting it
from \eqref{eq:content-projector} gives
\[
 \one_{G\le C}-1=-\one_{G>C},
\]
which is \eqref{eq:content-nontrivial-reassembly}.
\end{proof}

\subsection{\texorpdfstring{Physical reassembly at \(r=0\)}
{Physical reassembly at r=0}}

Apply \cref{lem:content-projector} to
\[
 G=c(\omega)
 =\gcd(N_\alpha(j),N_\gamma(j))
\]
for every literal atom \(\omega=(\alpha,\gamma,j)\).

\begin{theorem}[No content-projector centering]
\label{thm:content-no-centering}
At the physical cutoff \(C=\lfloor J\rfloor\),
\begin{equation}
 \boxed{
 Z_X
 =\sum_{\omega\in\cI_X}w_X(\omega)
  \sum_{t\mid c(\omega)}\eta_{\le C}(t)
 =\cS_{\rm full}^{\rm sh}-\cS_{>C}^{\rm sh}.}
 \label{eq:content-physical-reassembly}
\end{equation}
More precisely:
\begin{align}
 \text{\(t=1\) contribution}
 &=\cS_{\rm full}^{\rm sh},
 \label{eq:content-t-one}\\
 \text{sum of all \(t>1\) contributions}
 &=-\cS_{>C}^{\rm sh}.
 \label{eq:content-t-greater-one}
\end{align}
Thus the content projector does not annihilate or center the
distinguished zero mode.
\end{theorem}

\begin{proof}
Multiply \eqref{eq:content-projector} by \(w_X(\omega)\) and sum over
\(\omega\).  The left side is the small-content sum, hence \(Z_X\).
Because \(\eta_{\le C}(1)=1\), the \(t=1\) term sums every literal atom
with its original signed coefficient, which is
\(\cS_{\rm full}^{\rm sh}\).  Equation
\eqref{eq:content-nontrivial-reassembly} shows pointwise that the
remaining divisor terms equal \(-\one_{c(\omega)>C}\).  Their sum is
\(-\cS_{>C}^{\rm sh}\).
\end{proof}

\begin{corollary}[The principal term contains the hard problem]
\label{cor:content-hard-principal}
Using \eqref{eq:literal-large-content-bound},
\begin{equation}
 Z_X
 =\cS_{\rm full}^{\rm sh}+O(X^{o(1)}Q^2).
 \label{eq:content-full-mod-error}
\end{equation}
Therefore a bound
\(|Z_X|\ll X^{o(1)}Q^2\) is, at the same scale, a cancellation
statement for the unrestricted full matched shell.  It is not
obtained by estimating only the nonprincipal projector terms.
\end{corollary}

\begin{proof}
Combine \cref{thm:content-no-centering} with
\eqref{eq:literal-large-content-bound}.
\end{proof}

\begin{remark}[Exact reassembly versus termwise absolute bounds]
The useful estimate for the \(t>1\) terms comes from their exact
reassembly into the large-content shell.  Replacing that signed sum by
\[
 \sum_{\omega}|w_X(\omega)|
 \sum_{\substack{t\mid c(\omega)\\t>1}}
 |\eta_{\le C}(t)|
\]
can be much larger and does not create a bound for the \(t=1\) term.
The fact that nontrivial divisors appear in the expansion is not, by
itself, oscillation of the physical zero mode.
\end{remark}

\begin{remark}[Row calibration is a different operation]
The calibrated row prefix subtracts a row-dependent term and produces
two explicit drift legs; see \cref{rem:literal-calibration}.  Content
inversion acts on the common target gcd.  Neither identity says that
\(\sum_nA_C(n)\) vanishes, and the two operations must not be combined
into an unrecorded ``centered'' coefficient.
\end{remark}
